\section{Institutional Background}\label{sec:background}

\paragraph{The rule.} Article 11.8 of the collective bargaining agreement (CBA) between the NHL and the NHLPA allows the contract of any player who is a Group 3 unrestricted free agent to ``contain a no-Trade or a no-move clause'' \citep{cba2013}. A Group 3 free agent is a player with seven accrued seasons or aged 27 or older. A clause may be negotiated earlier only if it does not take effect until the player reaches that status. If the player is traded before the clause takes effect, the acquiring team is not bound unless it agrees in writing. Once in force, the clause stays in the contract after a trade. An NMC is broader than an NTC: it prevents ``the involuntary relocation of a Player, whether by Trade, Loan or Waiver claim,'' but it cannot restrict buyouts (Article 11.8(c)).

\paragraph{Contract practice.} The CBA does not define partial clauses, so the parties have invented them. A \emph{modified} NTC typically lets the player submit a list of teams to which he may not be traded (a no-trade list), or a list of teams he would accept (a trade list), often with a list size that changes over the life of the contract. A typical modern star contract combines several layers: an NMC for the first years, converting on a fixed date to a modified NTC whose list shrinks as the contract ages. Adam Lowry's Winnipeg extension, for instance, carries an NMC through June 2028 and then converts to a modified NTC with a five-team, and later a ten-team, trade list (CapWages contract notes). The median modified list in the data blocks ten teams. In the paper's notation, the player may block $k\approx10$ of the 31 other teams.

\paragraph{Term limits and the cap.} The 2005 CBA introduced the salary cap and set no maximum contract term, which produced a wave of very long, front-loaded deals. The 2013 CBA capped terms at eight years for a team re-signing its own player and seven otherwise. The 2025 agreement, effective 16 September 2026, cuts these to seven and six years and tightens salary variance \citep{dfo_dfo_cba}. The 2026 free-agency period was therefore the last under the old maximums, and teams hurried to register eight-year deals before the deadline. The cap was flat at \$81.5 million from 2019--20 through 2021--22 because of COVID-19. It has since risen sharply, to \$95.5 million in 2025--26 and \$104 million in 2026--27.

\paragraph{Why NMCs bind beyond trades.} For the 2017 Vegas and 2021 Seattle expansion drafts, players with an NMC in effect ``must be protected (and will be counted toward their club's applicable protection limits)'' \citep{nhl_expansion_rules}. Clauses therefore have real costs beyond trade demands, which the model captures as $C_0(k)$.

\paragraph{The comparison.} No-trade rights are far rarer in other North American leagues. NBA contracts almost never include them; one recent count puts the number of NBA players ever to hold one at about ten \citep{gerstle2025}. Major League Baseball has a statutory ``10-and-5'' right and occasional negotiated clauses. \citet{pedacehall2012} find that MLB players trade some salary for no-trade protection. In the NBA, star trade demands are frequent, but the player has no formal veto, so the resulting leverage is informal and comes mainly through threats not to re-sign. That is exactly the mechanism of Proposition~\ref{prop:informal}. Hockey is the league in which formal veto rights are most widespread, which makes it the natural laboratory.

\paragraph{Related work.} The economics of player mobility is old \citep{rottenberg1956,kahn2000,fortquirk1995}, and hockey salaries have been studied extensively \citep{idsonkahane2000,cyrenne2013,kahane2024}. \citet{fumarco2024} study strategic behaviour around NHL mobility restrictions, and \citet{marburger2009} asks why player trades dominate cash sales. As far as I can find, however, no peer-reviewed study estimates the price of NTCs or NMCs in the NHL, the league where they are most common, and none models trade demands as a strategic credibility problem. The closest descriptive work is \citet{hockeygraphs2021}. This paper aims to fill that gap.
